
%\needspace{3\baselineskip}
\color{uniblue}{\subsection{Контроль цепей напряжения стороны НН трансформатора (КЦН НН)}}
\color{black}

\begin{enumerate}[label=\arabic{section}.\arabic{subsection}.\arabic*, labelsep=4pt, leftmargin=0pt, itemindent=57pt]

\item
Контроль цепей напряжения стороны НН трансформатора служит для контроля исправности цепей напряжения стороны НН трансформатора.

\item
В устройстве контроль цепей напряжения представлен функцией «КЦН».

\item
Ввод функции в работу на этапе параметрирования устройства осуществляется программным переключателем SGF1 «Ввод\_функции». Информация о введенной функции отображается выходным сигналом «КЦН НН: Ввод».
\item
Функция может быть выведена из работы оперативно путем активации сигнала «ОВ КЦН НН», что характеризуется активным состоянием сигнала «КЦН НН: Оперативный вывод» на выходе функции.
\item
Функция контролирует снижение максимального из линейных напряжений ($U_{AB}, U_{BC}, U_{CA}$) ниже заданной уставки и повышение напряжения обратной последовательности ($U_2$) выше заданной уставки, в случае выполнения любого из условий через выдержку времени таймера Т1 (уставка «Тср») выдается сигнал неисправности ЦН.
\item
Алгоритм работы функции выполнен в соответствии с рисунком \ref{kznnn:img1}. В таблице \ref{kznnn:tbl1} приведены параметры, необходимые для настройки функции.

\vspace{3mm}
\begin{figure}[H]
\centering
\includegraphics[width=1\textwidth,height=1\textheight,keepaspectratio]{\fbpath/28. КЦН НН/_latex/img/1.pdf}
\captionof{figure}{Функциональная схема логики КЦН}\label{kznnn:img1}
\end{figure}

\end{enumerate}
\FloatBarrier % Все float'ы ДО этой строки будут размещены здесь.
